\section{Limitations and Conclusion}\label{sec:conclusion}

\odrift{} rests on two points: for OPE under drift, the age of a window matters less than
which mechanism changed it, and an estimator should rely on a drift decision only where the data
can observe that drift. A residual test on the conditional reward, unlike tests on raw
rewards, does not mistake context or logger drift for reward drift, and in our pre-registered
evaluation this made the difference between \odrift{} and a cross-domain estimator that groups
periods by mean reward in every reward-drift setting.

The limitations are as important as the result. First, drift-corrected reuse depends on the
reward-drift difference lying in the model class (A4); on semi-synthetic Open Bandit data,
with a misspecified model and large drift, it was heavily biased, and the estimator without
reuse was the better choice. Second, data-driven checks of reuse on the anchor or on a held-out
window could not repair this, because samples of the size available per period cannot detect
the biases involved; in our terms, the reuse bias is not observable at the anchor's sample size.
We see this as an information limit, not a tuning problem, and we expect
it to bind any method that validates reuse on the current period alone. Third, gradual drift
violates the piecewise assumption and favours methods that model a trend. Fourth, our
heterogeneity-adjusted interval undercovers when bias remains, and the certified alternative
is wide. Finally, Open Bandit provides real contexts and their real order but synthetic
rewards and drift; an evaluation on naturally drifting logged data remains open.

Two directions follow: reward models flexible enough to make (A4) plausible, so that reuse is
safe when attempted, and reporting for each estimate how much unobservable drift would change
a deployment decision.
